\begin{proposition}\label{prop:decic-no85}
An elliptic curve over a number field of degree $10$ has no point
of order $85$.
\end{proposition}
\begin{proof}
Put $X=X_1(85)$ and $J=J_1(85)$. Its genus is $225$.
The exact winding decomposition has two winding-zero constituents,
both of dimension $2$ and new level $85$. Their characters are
trivial on $2$. Every other constituent has nonzero winding.
By the winding rank-zero criterion of \cite[Section~3]{DEHMZ},
$(\langle2\rangle-1)J(\mathbb Q)$ is finite. We do not assume
that $J(\mathbb Q)$ has rank zero.
Let $D$ be the rational effective trace divisor of degree $10$
arising from a hypothetical level structure.

We first exclude entirely cuspidal reduction at $2$. We use the
multiplicative-level model and write $T_q^\mu,\delta_q$ for its
Hecke and diamond operators. On the two winding-zero constituents,
the characteristic polynomials at $2$ are
\[
 X^2+2X-1,\qquad X^2+2X+1,
\]
and at $3$ they are
\[
 X^2+4X+2,\qquad X^2+2X+2.
\]
Consequently the operators $k_2(T_2^\mu)$ and $k_3(T_3^\mu)$
have finite rational image, where
\[
 \begin{split}
 k_2(X)&=X^4+4X^3+4X^2-1,\\
 k_3(X)&=X^4+6X^3+12X^2+12X+4.
 \end{split}
\]
The exact polynomials and the two constituents are retained in the
accepted integral Hecke data.

Lift the geometric cusps of $\overline D$, with their multiplicities,
to a divisor $C$ over $\mathbb Q_2$. This is possible because the
cusps are finite etale away from $85$. The divisor $C$ need not be
rational over $\mathbb Q$. For a rational cusp $c$, each class
\[
 k_i(T_i^\mu)[D-C]
 =k_i(T_i^\mu)[D-10c]-k_i(T_i^\mu)[C-10c]
\]
is torsion and reduces to zero. The first term has finite rational
image; the second is torsion by Manin--Drinfeld.
In the multiplicative convention, the unramified dyadic formal-kernel
lemma gives the annihilator $T_2^\mu-1$. It is the point-level
annihilator $T_2-\langle2\rangle$ transported through Fricke and
multiplied by an invertible diamond. Thus
\[
 \Psi_i[D-C]=0,\qquad
 \Psi_i=(T_2^\mu-1)k_i(T_i^\mu),\quad i=2,3.
\]
We work throughout on the full Jacobian. No exactness of a quotient
cotangent sequence at $2$ is used.

The pullbacks of differentials by $\Psi_2,\Psi_3$ span a space $U$
of dimension $198$ over $\mathbb F_2$. In the saturated canonical
$q$-expansion basis the two binary matrices are
\[
 (T_2-I)(T_2^4+I),\qquad (T_2-I)T_3^4.
\]
Here $T_2$ is the canonical multiplicative-level matrix. In particular,
the factor is $T_2-I$, rather than $T_2-\delta_2$.
The space $U$ is stable under the full diamond group of order $32$.

The finite cusp signature is $1^{32}4^8 8^8$.
For reduction supported at rational cusps, translate a cusp of greatest
multiplicity to the first cusp. Retain ten differential jets there
and five at each other cusp. The resulting $165$ rows have rank
$152$. Independent enumeration of all $8191$ nonzero relations gives
minimum normalized cusp-prefix cost $32$. A relation occurring in a
divisor of degree $10$ would have cost at most $10$. Thus restriction
from $U$ has rank $10$ at every such rational-cuspidal divisor.

For the other cusps we compute genuine canonical arcs using
\[
 \begin{array}{c|c|c|c}
 \text{Tate curve}&\text{level point}&\text{residue degree}
       &\text{coefficients required}\\ \hline
 E(q^{17})&\zeta_5q&4&2\\
 E(q^5)&\zeta_{17}q&8&1\\
 E(q)&\zeta_{85}&8&1.
 \end{array}
\]
The global modular-unit identities recover these arcs from full-rank
canonical minors. A second computation uses different row minors and
recovers the same coefficients. Every retained unit identity is
checked, and the quartic arc has nonzero canonical tangent.
The exact diamond action gives the eight quartic places and the two
orbits of four octic places. Coefficientwise Frobenius on each arc
is checked against the appropriate diamond, so restriction of scalars
gives the correct closed-cusp conditions.

If a quartic cusp occurs, move one to the first quartic place.
The remaining patterns are one quartic and rational degree $6$, or
two quartics and rational degree $2$, allowing multiplicity $2$ at
the same quartic cusp. If an octic cusp occurs, it has multiplicity
$1$ and leaves rational degree $2$. The complete normalized counts are
\[
 \binom{37}{6}+8\binom{33}{2}+2\binom{33}{2}=2330064.
\]
Every restriction matrix has rank $10$. A separate implementation
checks each complete matrix, and independently verifies the rational
relation space and its minimum cost. Hence the product of the two
maps $\Psi_i$ is a formal immersion on $X^{(10)}$ at every cuspidal
reduction. Since $D$ and $C$ have the same reduction and image, we
obtain $D=C$, a contradiction.

It remains to exclude noncuspidal reduction. The complete independent
finite-field census gives no noncuspidal places below degree $9$,
then $160$ and $128$ in degrees $9$ and $10$ respectively.
Thus $\overline D$ has one degree-$9$ noncuspidal place and a rational
cusp, or is a degree-$10$ noncuspidal place.
In the point-level convention put
\[
 \sigma=\langle2\rangle^{-1},\qquad \tau=\langle3\rangle,
 \qquad A_3=T_3-1-3\tau.
\]
The class $[\sigma D-D]$ is rational torsion. The same holds for
the rational cuspidal complement. Odd-prime torsion injectivity and
Eichler--Shimura show that, for the noncuspidal place $P$ and $T=T_3P$,
\[
 B=\sigma T-T-\sigma P+P-3\sigma\tau P+3\tau P
\]
must be principal. The accepted degree-$9$ canonical certificate
proves the contrary for all five source orbits. Its only dependence
on the total degree is the rationality of the cuspidal complement,
which remains valid here.

The new degree-$10$ curves have $1020$ points and form four closed
$j$-orbits. Their groups are cyclic: the smaller invariant factor
must divide $1023$ and have square dividing $1020$. Thus their
$85$-torsion gives four free diamond orbits of length $32$, accounting
for all $128$ source points. For one representative of each orbit,
the actual four geometric $3$-kernels give target places of degrees
$10,30$, with norm multiplicities $1,1$. Independent Sage Velu
isogenies verify every actual point image and both norm factors.
The ordinary deformation $a_6\mapsto a_6+t$ and the $85$-division
recurrence give genuine canonical jets through order $5$.
All twelve source and target families have unit-recovery rank $224$
and nonzero tangent, checked independently.

The four required divisors $B$ have positive and negative degree $80$.
We use $R_1=H^0(X,K_X)$ and $R_2=H^0(X,2K_X)$, of dimensions
$225,672$, with their accepted exact multiplication table.
Choose nonzero $F,H\in R_1$ vanishing on $B_+$. If
$f\in\mathcal L(B)$, then $G=fF$ vanishes on $B_-$ and
$GH\in FR_1$. In each test multiplication by $F$ has rank $225$,
and the joint restriction and multiplication map
\[
 R_1\longrightarrow H^0(B_-,K_X|_{B_-})\oplus R_2/FR_1,
 \qquad G\longmapsto(G|_{B_-},GH\bmod FR_1)
\]
has rank $225$. An independent Sage calculation reconstructs the
quotient and gives the same rank. Hence $\mathcal L(B)=0$, a
contradiction.
\end{proof}
